%% ma: L12-B08-0059
%% lop: 12
%% chuong: 2
%% bai: 8
%% dang: 12.8.35
%% loai: DS
%% muc: [NB, NB, TH, VD]
%% dap_an: "ĐĐĐĐ"
%% nguon: "LaTeX giáo viên gửi 21/10/2026 (M05, Vectơ và hệ trục tọa độ trong không gian, 107 câu), Phần II Câu 14"
%% kien_thuc: "Bài 7 (tọa độ điểm), Bài 8 (tích vô hướng, điều kiện vuông góc)"
%% trang_thai: cho-duyet
%% ly_do: "Dạng mới đề xuất 12.8.35; Đã bổ sung cách chọn trục Ox, Oy (đề gốc thiếu nên ý a) không xác định được), cần giáo viên xác nhận"
%% ghi_chu: "Đề gốc chỉ nói 'trục Oz chứa đường cao OS', chưa xác định trục Ox, Oy nên ý a) không xác định được; đã bổ sung 'các tia Ox, Oy, Oz lần lượt chứa các tia OA, OB, OS'"
\begin{ex}
Cho hình chóp $S.ABCD$ có đáy $ABCD$ là hình vuông tâm $O$, $OA=OB=OC=OD=a$, $SO$ vuông góc với mặt đáy và $SO=h$. Chọn hệ trục tọa độ $Oxyz$ sao cho các tia $Ox$, $Oy$, $Oz$ lần lượt chứa các tia $OA$, $OB$, $OS$.
\choiceTF
{\True Tọa độ các đỉnh của mặt đáy là $A(a;0;0)$, $B(0;a;0)$, $C(-a;0;0)$, $D(0;-a;0)$.}
{\True Tọa độ đỉnh $S$ là $(0;0;h)$.}
{\True Tích vô hướng $\overrightarrow{SA}\cdot\overrightarrow{SC}=h^2-a^2$.}
{\True Nếu $\overrightarrow{SA}\perp\overrightarrow{SC}$ thì chiều cao $h$ bằng độ dài đoạn $OA$.}
\loigiai{a) $OA\perp OB$ (hai đường chéo hình vuông), $C$, $D$ đối xứng với $A$, $B$ qua $O$. Đúng.

b) $S$ thuộc tia $Oz$, $OS=h$. Đúng.

c) $\overrightarrow{SA}=(a;0;-h)$, $\overrightarrow{SC}=(-a;0;-h)$, $\overrightarrow{SA}\cdot\overrightarrow{SC}=h^2-a^2$. Đúng.

d) $\overrightarrow{SA}\perp\overrightarrow{SC}\Leftrightarrow h^2=a^2\Leftrightarrow h=a=OA$. Đúng.}
\end{ex}
